
\begin{lemma}{7.2.336}\label{lem:7-2-336}
对满足上述等式
(7.2.335.1)、(7.2.335.2)
和 (7.2.335.3) 的任意
$\mathcal F_d^{s=s_0}$，
存在滤过 $A_{n',K'}$ 同态
\[
\mathfrak H_d\colon
C(L_1,(\lambda_0,\lambda_1))
\longrightarrow
\widehat{\mathfrak C}_d(L'_1,\lambda'_1),
\]
满足下列性质：
\begin{equation}
\operatorname{Eval}_{r=0}\circ\mathfrak H_d
=\mathfrak e_d\circ\mathfrak E_d
\circ\mathfrak f_{\mathrm{ext}}^{\{\psi_{1,s}\}_s},
\tag{7.2.337.1}\label{eq:7-2-337-1}
\end{equation}
\begin{equation}
\operatorname{Eval}_{r=1}\circ\mathfrak H_d
=\mathfrak e_d\circ
\mathfrak f^{\{\psi_{1,s}\}_s}\circ j,
\tag{7.2.337.2}\label{eq:7-2-337-2}
\end{equation}
\begin{equation}
\operatorname{Eval}_{s=s_0}\circ\mathfrak H_d
=\mathcal F_d^{s=s_0},
\tag{7.2.337.3}\label{eq:7-2-337-3}
\end{equation}
\begin{equation}
\mathfrak H_d
=\operatorname{Incl}_r\circ\mathfrak e_d
\circ\mathfrak f^{\{\psi_{1,s}\}_s}
\quad\text{在 }C(L_1,\lambda_1)\text{ 上}.
\tag{7.2.337.4}\label{eq:7-2-337-4}
\end{equation}
\end{lemma}

\noindent\textbf{证明。}
应用命题~\ref{prop:7-2-312}。
取
\[
C'_0=C(L_1,\lambda_1),\qquad
C'=C(L_1,(\lambda_0,\lambda_1)),\qquad
\widehat{\mathfrak C}
=\widehat{\mathfrak C}_d(L'_1,\lambda'_1).
\]
令
$\mathfrak H^{r=0}$、
$\mathfrak H^{r=1}$、
$\mathfrak H^{s=s_0}$
分别为
(7.2.337.1)、(7.2.337.2)、
(7.2.337.3)
的右端。
又取
\[
\mathfrak H_0
=\operatorname{Incl}_r\circ\mathfrak e_d
\circ\mathfrak f^{\{\psi_{1,s}\}_s}.
\]
令
\[
\mathfrak h\colon
C(L_1,\lambda_1)
\longrightarrow C(L_1,(\lambda_0,\lambda_1))
\]
为包含映射。

下面验证
(7.2.313.1)、(7.2.313.2)
和 (7.2.313.3)。
由 (7.2.333.4)、(7.2.333.5)、
(7.2.331.2) 及 (7.2.335.1)，得到
\[
\begin{aligned}
\operatorname{Eval}_{s=s_0}\circ\mathfrak e_d
\circ\mathfrak E_d
\circ\mathfrak f_{\mathrm{ext}}^{\{\psi_{1,s}\}_s}
&=j_{s=s_0}\circ j_{s=s_0}
\circ\mathfrak f_{\mathrm{ext},s=s_0}^{\psi_{1,s_0}}\\
&=\operatorname{Eval}_{r=0}
\circ\mathcal F_d^{s=s_0}.
\end{aligned}
\]
同样由 (7.2.333.5)、(7.2.331.2)
及 (7.2.335.2)，计算得
\[
\begin{aligned}
\operatorname{Eval}_{s=s_0}\circ\mathfrak e_d
\circ\mathfrak f^{\{\psi_{1,s}\}_s}\circ j
&=j_{s=s_0}\circ
\mathfrak f^{\psi_{1,s_0}}\circ j\\
&=\operatorname{Eval}_{r=1}
\circ\mathcal F_d^{s=s_0}.
\end{aligned}
\]
因此 (7.2.313.1) 得证。

为验证 (7.2.313.2)，
由图式~\ref{diag:7-2-23} 的严格交换性
与 (7.2.333.3)\translatornote{原书这里引用
(7.2.333.1) 与 (7.2.333.3)；
但 (7.2.333.1) 只给出同伦交换，
不能直接推出所需的严格等式。
图式 7.2.23 的严格交换性
及 $\mathfrak E_d$ 在小模型上的恒等性
足以推出该等式。}，
在 $C(L_1,\lambda_1)$ 上有
\[
\begin{aligned}
\mathfrak e_d\circ\mathfrak E_d
\circ\mathfrak f_{\mathrm{ext}}^{\{\psi_{1,s}\}_s}
&=\mathfrak e_d
\circ\mathfrak f^{\{\psi_{1,s}\}_s}\\
&=\operatorname{Eval}_{r=0}\circ
\operatorname{Incl}_r\circ\mathfrak e_d
\circ\mathfrak f^{\{\psi_{1,s}\}_s}.
\end{aligned}
\]
根据 $j$ 的定义，
在 $C(L_1,\lambda_1)$ 上还成立
\[
\begin{aligned}
\mathfrak e_d\circ
\mathfrak f^{\{\psi_{1,s}\}_s}\circ j
&=\mathfrak e_d\circ
\mathfrak f^{\{\psi_{1,s}\}_s}\\
&=\operatorname{Eval}_{r=1}\circ
\operatorname{Incl}_r\circ\mathfrak e_d
\circ\mathfrak f^{\{\psi_{1,s}\}_s}.
\end{aligned}
\]
故 (7.2.313.2) 得证。

最后验证 (7.2.313.3)。
由 (7.2.335.3)，
在 $C(L_1,\lambda_1)$ 上有
\[
\mathcal F_d^{s=s_0}
=\operatorname{Incl}_r\circ j_{s=s_0}
\circ\mathfrak f^{\psi_{1,s_0}}.
\]
由 (7.2.334)、(7.2.333.5)
及 (7.2.331.2)，
在 $C(L_1,\lambda_1)$ 上又有
\[
\begin{aligned}
&\operatorname{Eval}_{s=s_0}\circ
\operatorname{Incl}_r\circ\mathfrak e_d
\circ\mathfrak f^{\{\psi_{1,s}\}_s}\\
&\qquad
=\operatorname{Incl}_r\circ
\operatorname{Eval}_{s=s_0}\circ\mathfrak e_d
\circ\mathfrak f^{\{\psi_{1,s}\}_s}\\
&\qquad
=\operatorname{Incl}_r\circ j_{s=s_0}
\circ\mathfrak f^{\psi_{1,s_0}}.
\end{aligned}
\]
这就证明了 (7.2.313.3)。
引理~\ref{lem:7-2-336}
于是由命题~\ref{prop:7-2-312}
推出。\hfill$\square$
